%%%%%%%%%%%%%%%%
% PROBLEMA
%%%%%%%%%%%%%%%%

\section{Planteamiento del Problema}

En los procesos actuales de transformación digital, las organizaciones han aumentado la demanda de desarrollo y evaluación de software con el fin de soportar sus capacidades de negocio. Esta tendencia se enmarca en el reconocimiento de la transformación digital como un impulsor fundamental de la competitividad empresarial, con cerca del 90\% de las organizaciones inmersas en estas iniciativas y un gasto global proyectado que superará los 4 billones de dólares para 2027 \citep{Jie2025}. Este crecimiento ha generado una presión constante por reducir los tiempos de entrega y aumentar la eficiencia de los equipos de desarrollo, debido a que la velocidad de salida al mercado se ha convertido en un factor diferencial para la competitividad y la satisfacción de los usuarios.

Como respuesta a estos desafíos, las metodologías ágiles se han convertido en un enfoque clave para la gestión y el desarrollo de proyectos \citep{Palopak2024}. Sin embargo, en la práctica moderna de la ingeniería de software, especialmente en entornos ágiles, los equipos enfrentan diversas tensiones de desempeño, como el equilibrio entre velocidad y calidad, lo que evidencia la persistencia de una brecha entre la velocidad esperada por el negocio y la capacidad real de los equipos para mantener dicha velocidad con calidad técnica \citep{Armbruster2025}.

Evidencias de esta brecha se encuentran en casos como el Laboratorio Digital del Banco W \citep{UniversidadIcesi2023}, el cual opera bajo marcos de trabajo ágiles, incluyendo prácticas como Dual Track Agile \citep{MartyCagan2012}. Este modelo, aunque mejora la velocidad de innovación y la capacidad de respuesta, también genera un desafío. A medida que el equipo logra entregar más rápido, las expectativas institucionales aumentan, transformando la agilidad de una ventaja a un nuevo estándar que exige entregas aún más rápidas y frecuentes.

A pesar de esto, en la práctica habitual de desarrollo en el Laboratorio Digital, el inicio de nuevos proyectos implica configurar desde cero elementos como la estructura base del sistema, las convenciones técnicas y las integraciónes con servicios externos, concentrando tiempo en actividades que no generan valor agregado al producto. De manera similar, la incorporación de funcionalidades transversales entre aplicaciones, como autenticación, observabilidad o mensajería, se realiza frecuentemente a través de procesos manuales, reutilizando código existente, dado que una parte significativa de la base técnica tiende a repetirse a lo largo de diferentes aplicaciones. Como resultado, se generan duplicidad de funcionalidades e integraciones inconsistencias en la estructura de los proyectos que dificultan su mantenimiento y la evolución controlada del ecosistema de aplicaciones en la organización.

Bajo esta situación, existen soluciones que buscan encapsular la estructura y organización de los proyectos mediante plantillas base o boilerplates, entendido como conjuntos de código y configuraciones predefinidas que permiten iniciar nuevos desarrollos sobre una base común. No obstante, su reutilización real suele ser limitada, ya que estos mecanismos tienden a materializarse como artefactos estáticos que rápidamente quedan desactualizados o se acoplan a contextos específicos. Incluso, se ha señalado que los boilerplates rara vez logran ser reutilizables de forma general debido a la diversidad de necesidades y tecnologías entre proyectos \citep{SeanPromsopeak2022}.

En consecuencia, persiste la necesidad de establecer una base técnica robusta que le permita a equipos de desarrollo, como el Laboratorio Digital del Banco W, industrializar el arranque de proyectos y la composición de capacidades reutilizables sin sacrificar la autonomía de los equipos ni los atributos de calidad del software.